% !TeX root = ../main.tex

\chapter{绪论}
\label{cha:introduction}

\section{研究背景与意义}
\label{sec:intro-background}

自动驾驶系统的研究重点正在从常规场景中的目标检测、跟踪与轨迹预测，逐步转向低频、高风险且语义开放的长尾场景理解。对于真实道路而言，影响系统安全边界的往往并不是对常见目标的普遍失效，而是对远距离目标、遮挡边界、局部关键证据以及复杂上下文关系的理解不足\cite{caesar2020nuscenes,dosovitskiy2017carla,liu2024surveyadatasets,kaur2021surveyselfdrivingsimulators}。在行人突然出现、施工区域临时改道、异常交通参与者穿行等场景中，系统不仅要判断“看见了什么”，还要回答“场景意味着什么”“应当如何响应”，这使得自动驾驶感知问题逐步从封闭类别识别扩展为兼具语义理解和风险解释的综合任务。

视觉语言模型（Vision-Language Model, VLM）通过统一视觉输入与语言输出，为自动驾驶场景理解提供了新的研究方向。相较于传统封闭类别感知模型，视觉语言模型在开放语义问答、多对象关系表达和任务条件化理解等方面表现出更强的表达能力，这使其在驾驶场景问答、场景摘要、异常解释和人机交互辅助等任务中展现出显著潜力\cite{marcu2024lingoqa,sima2024drivelm,guo2024drivemllm,wang2024omnidrive}。从研究趋势看，将视觉语言模型引入自动驾驶，不仅是模型规模演进的结果，更反映了自动驾驶系统对“可解释场景语义”的现实需求。

然而，视觉语言模型在自动驾驶中的部署仍面临鲜明的系统约束。车端边缘侧需要低时延、高稳定性、可持续运行与受控能耗，而长尾场景理解又常常需要更强的视觉语义推理能力。若完全依赖大模型常驻推理，系统难以满足实时性与成本预算；若完全依赖边缘小模型，则又难以稳定覆盖远距离目标和低频异常情形。与此同时，传统通过数据回流、离线标注与重新训练来修复长尾失效的路径，虽然在模型研究中常见，但在实际系统中往往存在反馈周期长、更新成本高和新场景响应慢等问题。因而，如何在冻结主模型权重的前提下，以系统级组织方式实现长尾场景下的快速补救与经验积累，已成为面向部署的自动驾驶感知研究中的重要问题。

\section{研究问题定义与部署约束}
\label{sec:intro-problem}

围绕上述背景，本文聚焦如下核心问题：在不更新主模型参数的条件下，能否构建一种同时兼顾边缘实时性、云端纠错能力与知识复用能力的自动驾驶长尾场景感知框架？该问题本质上属于“部署约束下的测试时增强”问题，其关注重点不是进一步提升离线训练精度，而是在推理阶段通过角色分层、路径分层与外部知识组织获得系统增益。

更具体地，本文将问题拆解为三个相互关联的子问题。

\begin{enumerate}
  \item 边缘端小型视觉语言模型在距离敏感任务中的性能退化，能否通过系统级边缘云协同获得稳定补救？
  \item 云端能力在系统中应当扮演何种角色，即它究竟应被理解为默认主链路的替代方案，还是应被组织为受约束的低频补偿机制？
  \item 在冻结权重的前提下，单次纠错经验能否被抽象为可检索、可约束、可复用的外部知识，从而形成测试时自适应能力？
\end{enumerate}

为了使该问题具有明确的研究边界，本文进一步强调四项部署约束。第一，系统输出不能停留在自由文本答案层面，而应具有可审计、可分析的结构化中间表示。第二，云端调用必须受到代价约束，不能以牺牲实时性为代价换取局部准确率提升。第三，知识复用机制必须具备风险边界，避免错误经验在相似场景中被无约束扩散。第四，实验评测既要关注任务效果，也要同时覆盖时延、云调用比例、错误修复率和失效模式解释等系统指标。

\section{现有研究不足与研究空白}
\label{sec:intro-gap}

现有相关研究已经分别从自动驾驶场景问答、多模态大模型、边缘云协同和测试时自适应等角度开展探索，但在本文关注的研究交叉点上仍存在明显空白。

首先，自动驾驶多模态研究多数围绕离线基准表现展开，对部署场景中的边缘算力、网络波动和云调用预算关注不足。已有工作通常默认更强模型始终可用，因此更关注“模型能否答对”，而较少回答“系统在受限预算下如何答对”这一部署导向问题。其次，现有视觉语言模型在远距离目标、小尺度区域和局部关键证据利用方面仍存在明显短板，尤其在驾驶问答任务中容易出现假阴性偏置、过度保守判断和证据锚定错误，但现有研究对这些失效模式与系统补救路径之间的对应关系阐述不足。

再次，测试时增强研究已经表明，能力提升并不一定只能通过参数更新获得，还可以通过工具调用、检索增强和外部记忆实现\cite{yao2023react,shinn2023reflexion,madaan2023selfrefine,schick2023toolformer,lewis2020rag,guu2020realm}。然而，这些研究大多面向通用文本任务或开放式智能体环境，对自动驾驶场景中特有的低时延要求、结构化输出需求和安全回退机制讨论不足。换言之，现有工作已经说明“外部知识可能有用”，但尚未充分回答“在自动驾驶中怎样的外部知识表示是可审计、可约束且可复用的”。

最后，从实验叙事角度看，系统研究容易出现证据层级混杂的问题。例如，将在线路由收益与长期知识沉淀收益写在同一组主实验中，容易导致结论边界模糊；再如，将小规模适配原型中的局部正向结果直接提升为普遍结论，也会削弱论文的可信度。因此，如何建立与证据强度相一致的研究叙事，同样是当前研究中仍待完善的关键环节。

\section{研究动机与总体思路}
\label{sec:intro-motivation}

上述研究空白共同指向一个核心动机：自动驾驶长尾场景中的有效增强，不应仅被理解为“接入更大模型”，而应被理解为“在受约束的系统链路中，将不同能力组织为互补关系”。基于这一判断，本文不以参数更新为主要路径，而是将边缘模型、云端模型与外部知识视为三类功能互补的能力来源。

其中，边缘模型负责高频、低时延的基础场景感知；云端模型负责在关键脆弱样本上提供低频但高价值的纠错；结构化技能库则负责将纠错经验转化为外部知识单元，为后续相似场景提供受约束的再感知指导。由此，本文形成“边缘感知、不确定性评估、选择性云协同、结构化技能沉淀与约束化再感知”的闭环机制。该思路的基本出发点在于：长尾场景中的失败并非均匀分布，而是集中出现在若干高风险条件下，因此更适合通过选择性系统补救而非无差别全量重算来处理。

进一步地，本文将研究目标划分为两个层面。第一层面是在线系统有效性，即在不引入大规模云调用的前提下，选择性混合路由是否能够补救关键长尾样本。第二层面是长期经验外化，即单次纠错经验能否被整理为结构化知识，并在后续样本中形成稳定的正向复用。前者回答“系统是否有效”，后者回答“经验能否沉淀”，两者在研究目标上相互关联，但在证据层级上必须严格区分。

\section{研究思路与主要贡献}
\label{sec:intro-contribution}

针对上述问题，本文提出一种免训练的分层边缘云协同感知框架。该框架以边缘端小型视觉语言模型承担高频主链路，以云端多模态模型承担低频纠错链路，并通过结构化技能沉淀机制将纠错经验转化为外部知识单元。与单纯追求更大模型或更高离线基准分数的路径不同，本文关注的是在冻结权重条件下探索更贴近部署需求的系统方法。

本文的主要贡献概括如下。

\begin{enumerate}
  \item 从部署约束出发，提出自动驾驶长尾场景中的免训练分层边缘云协同感知问题，明确给出输入输出、代价约束、角色分层和证据边界，并建立边缘感知、云端纠错与外部知识沉淀相统一的研究框架。
  \item 设计结构化技能沉淀机制，将单次纠错经验表示为包含触发条件、关注区域、问题分解、输出约束与安全回退的外部知识单元，从而为测试时知识外化与受约束复用提供可操作的表示形式。
  \item 建立“基准忠实协议”和“适配评估协议”相分离的实验叙事，在 DTPQA 清洗共享子集上验证选择性混合路由的系统有效性，并通过独立适配实验与真实数据试验分析方法的机制边界、风险来源与外部效度。
\end{enumerate}

\section{论文结构安排}
\label{sec:intro-organization}

全文共分为六章，各章内容安排如下。

\begin{enumerate}
  \item 第一章为绪论，介绍研究背景、核心问题、主要挑战与本文贡献。
  \item 第二章为相关工作，综述自动驾驶场景理解基准、驾驶多模态大模型、视觉语言模型感知短板以及测试时自适应研究。
  \item 第三章给出问题定义与总体框架，形式化描述输入输出、优化目标、系统结构与证据边界。
  \item 第四章介绍本文方法，重点说明边缘感知、路由决策、云端纠错、结构化技能沉淀与约束化再感知机制。
  \item 第五章报告实验设置、主结果、机制分析与真实数据试验，讨论方法的有效性与适用边界。
  \item 第六章总结全文，并从部署意义、方法局限与未来研究方向三个方面展开讨论。
\end{enumerate}
